\documentclass[../../../include/open-logic-section]{subfiles}

\begin{document}
\olfileid{sth}{ordinals}{zfm}
\olsection{$\ZFminus$: tonggak wigati}

Pitakonan carane ndhukung Panggantian (yen pancen bisa)
ora gampang. Mula kita bakal ngrembug ing
\olref[replacement][]{chap}. Nanging kanthi nambah
Panggantian, kita wis tekan tonggak wigati liyane.
Saiki kita nduweni kabeh aksioma kang dibutuhake
kanggo teori $\ZFminus$. Rincine mangkene:

\begin{defn}
Teori $\ZFminus$ nduweni aksioma-aksioma iki:
Ekstensionalitas, Gabungan, Pasangan, Himpunan Pangkat,
Tanpa Wates, lan kabeh conto saka skema Pamisahan lan
Panggantian. Kanthi tembung liyane, $\ZFminus$
nambahake Panggantian marang $\Zminus$.
\end{defn}
\noindent
Jeneng iki nuduhake teori himpunan
\emph{Zermelo--Fraenkel} (\emph{dikurangi} sawijining
bab kang bakal kita rembug mengko). Fraenkel diajeni
amarga dheweke diakoni ngrumusake Panggantian ing
\citeyear{Fraenkel1922}, sanadyan rumusan trep kang
kapisan digawe dening \citet{Skolem1922}.

\end{document}
